\section{Pre-registration (verbatim)}
\label{app:prereg}
This section reproduces the analysis plan as frozen before any AIPR score was
joined to any outcome. The primary hypothesis (H1), the primary metric
(bottom-quintile reject lift), the discrimination metric (AUROC), the trend test
(Jonckheere--Terpstra, Monte Carlo permutation), the bridge gate ($\rho\ge0.8$),
and the exclusion rules were all fixed in advance. The text below is reproduced
verbatim from \texttt{DECISIONS.md} at the public tag
\texttt{prereg-iclr2026-v1} (\texttt{github.com/cogeor/aipr-score-validation});
the push/tag timestamp on that public remote is the freeze anchor (commit dates
are forgeable). The verbatim block is generated from the committed file by
\texttt{cli.py paper}, so it cannot drift from the frozen pre-registration.

\input{prereg_verbatim}

\section{Grading configurations}
\label{app:configs}
\begin{table}[h]
  \centering
  \caption{Grading configurations. The two pipeline configurations are the cost
  ladder (cohort nesting $H\subseteq M$); they run the identical full two-pass
  pipeline and share the scoring formula (Eq.~\ref{eq:overall}), differing only in
  the reviewer model tier. The naive judge (Appendix~\ref{app:naive}) is graded on
  cohort $H$ so it pairs with a full-text score; it uses the same frontier model
  but a single prompt with no rubric or audit, and returns only an overall score.}
  \label{tab:configs-app}
  \begin{tabular}{lll}
    \toprule
    Configuration & Pipeline input / mode & Model tier \\
    \midrule
    full-mini  & full text, two-pass (+ audit)     & cheap \\
    full       & full text, two-pass (+ audit)     & frontier \\
    \midrule
    naive      & full text, single one-paragraph prompt (no rubric/audit) & frontier \\
    \bottomrule
  \end{tabular}
\end{table}

\section{Data schema}
\label{app:schema}
The analysis consumes two tables. \texttt{submissions} (one row per submission):
identifier, venue, year, decision tier, tier rank, accept flag, mean reviewer
rating, rating SD, review count, sampling stratum, an exclusion flag with reason,
and optional manuscript metadata (primary area, page/word/token/reference counts).
\texttt{gradings} (one row per grading run): submission id, configuration, run
index, model name, pipeline version, the five dimension scores, the overall score,
the tokens used, and provenance ids. Contract invariants (tier/rank consistency,
uniform pipeline version, score ranges, referential integrity, cohort nesting
$H\subseteq M$) are asserted at load.

\section{Design justification: statistical power and estimator validity}
\label{app:power}
Both checks below were run \emph{before} data collection and do not use the real
outcomes; they establish that the planned sample is adequate and that the
analysis code is correct. They are reproduced by \texttt{analysis/simulation.py}.

\paragraph{Power.}
Under a generative model in which a latent paper quality drives the decision
tier, the reviewer rating, and (more noisily) the AIPR score, we swept the
effect-size knob and report power as a function of the \emph{true} AUROC rather
than a single guessed value (Fig.~\ref{fig:power}). At the planned cohort sizes
the design is well powered: for a moderate true effect the frontier cohort
($n=\NfullPrimary$) attains power $\powerHtwoN\%$ on the separation test (H2) and
$\powerHoneN\%$ on the primary low-end test (H1), and the full-mini cohort
($n\approx\NminiPrimary$) is at ceiling. The minimum detectable effect for H2 at
$n=100$ and 80\% power is a true AUROC of about $\mdeAUROC$, well below any
plausible value, so a null result would be informative rather than
underpowered.

\paragraph{Estimator validity.}
On data with a known population value, our 95\% BCa bootstrap confidence intervals
covered the truth $\bootCoverage\%$ of the time (nominal 95\%; the simpler
percentile interval under-covers at this sample size, which is why we use BCa). Under a true null
(no trend across tiers), the Jonckheere--Terpstra permutation test rejected at
$\jtTypeOne\%$ (nominal 5\%) and its $p$-values were uniform (Kolmogorov--Smirnov
$p=\jtUniformP$). The analysis therefore controls error at its stated levels.
Additionally, a label-shuffle null control is asserted at analysis time: with
decision labels permuted, the overall-score AUROC is $\approx 0.5$.

\begin{figure}[h]
  \centering
  \includegraphics[width=\textwidth]{figS_power}
  \caption{Pre-data power analysis. Left: power of the separation test (H2) vs.\
  true effect size at the two planned cohort sizes; dotted line is 80\% power.
  Right: power of all four hypotheses at the frontier cohort ($n=100$). Both
  generated by \texttt{simulation.py} on a synthetic generative model, not on
  study data.}
  \label{fig:power}
\end{figure}

\section{Robustness and secondary analyses}
\label{app:robust}

\begin{figure}[h]
  \centering
  \includegraphics[width=0.7\columnwidth]{figS_roc}
  \caption{ROC for reject-vs-accept discrimination, full-mini (headline) and frontier
  (production) configurations; legends report AUROC with 95\% CIs (H2). The AUROC
  scalars also appear in Table~\ref{tab:headline} and the deployable-end
  discrimination is the reliability panel Fig.~\ref{fig:validation}b.}
  \label{fig:roc}
\end{figure}

\begin{figure}[h]
  \centering
  \includegraphics[width=0.75\columnwidth]{figS_subscore_auroc}
  \caption{Per-dimension reject-vs-accept AUROC (full-mini cohort) with 95\% CIs.
  Novelty and applicability carry the most discriminative signal; clarity the
  least.}
  \label{fig:subscore-auroc}
\end{figure}

\begin{table}[h]
  \centering
  \caption{Per-dimension AUROC (full-mini cohort).}
  \label{tab:subscores-app}
  \input{tables/tab_subscores}
\end{table}

\begin{figure}[h]
  \centering
  \begin{subfigure}{0.49\columnwidth}
    \includegraphics[width=\linewidth]{figS_nested_auroc}
    \caption{Cheap vs.\ frontier model on cohort $H$.}
  \end{subfigure}\hfill
  \begin{subfigure}{0.49\columnwidth}
    \includegraphics[width=\linewidth]{figS_run_variance}
    \caption{Run-to-run scoring variance.}
  \end{subfigure}
  \caption{Left: AUROC for full-mini and full on the shared cohort $H$, isolating
  whether the frontier model adds discriminative signal over the cheap model (both
  run the full two-pass pipeline). Right: distribution of within-paper score SD
  across repeated frontier runs.}
  \label{fig:supp-pair}
\end{figure}

\paragraph{Multiple comparisons.} The five per-dimension discrimination tests are
controlled at FDR $0.05$ by the Benjamini--Hochberg procedure; adjusted $q$-values
accompany Table~\ref{tab:subscores-app}. The confirmatory hypotheses H1--H4 are
pre-specified single tests and are not part of this family.

\begin{figure}[h]
  \centering
  \includegraphics[width=0.6\columnwidth]{figS_prevalence}
  \caption{Low-score flag under prevalence shift. AUROC is prevalence-invariant,
  but a flag's precision is not---we re-weight the balanced cohort across a range
  of assumed natural acceptance rates and show the lowest-score band stays well
  above the base reject rate throughout.}
  \label{fig:supp-cond}
\end{figure}

\paragraph{Prevalence re-weighting.} Because cohort $M$ over-samples the rarer
accept tiers relative to the venue's natural prevalence,
we re-weight to a range of natural acceptance rates (Fig.~\ref{fig:supp-cond});
the discrimination (AUROC) is prevalence-invariant and the low-score
flag's precision remains above base rate across the plausible range. At the
venue's actual prevalence (\natAcceptRate\% accept, base reject rate
\natBaseReject\%), the bottom-quintile reject precision is \natBottomPrecision\%.

\paragraph{Second-venue replication.} The headline discrimination and rating
correlation are recomputed on a second venue/year (\repVenue) on the cheap-model
full-mini configuration; AUROC~\aurocRepFull (Fig.~\ref{fig:replication}).

\begin{figure}[h]
  \centering
  \includegraphics[width=0.6\columnwidth]{figS_replication}
  \caption{Reject-vs-accept ROC on the replication venue (\repVenue, full-mini
  configuration). Same direction and magnitude as the primary cohort,
  supporting year-to-year stability of the signal.}
  \label{fig:replication}
\end{figure}

\paragraph{Mini-to-frontier bridge (H5).} On the paired cohort $H$, the cheap-model
full-mini score and the frontier full score agree closely
(Fig.~\ref{fig:bridge}, left; $\rho =$~\bridgeRhoFull, $n=\bridgeN$), clearing the
pre-registered $\rho\ge0.8$ gate that licenses the large-$N$ full-mini results as a
proxy for the production score (\S\ref{sec:results}). Because a global
correlation can mask low-end disagreement, we also report agreement restricted to
the full-mini bottom quintile, $\rho =$~\bridgeLowRhoFull, with \bridgeOverlap\% of the
full-mini bottom quintile also in the frontier bottom quintile
(Fig.~\ref{fig:bridge}, right): the proxy holds where the deployable claim lives.

\begin{figure}[h]
  \centering
  \includegraphics[width=\textwidth]{fig5_bridge}
  \caption{Mini-to-frontier bridge on the paired cohort $H$. Left: full-mini vs.\ full
  overall score, dashed line identity (global agreement). Right: bottom-quintile
  membership --- papers in both scores' bottom quintile vs.\ mini-only; dotted
  lines are the per-score quintile thresholds. The low-end recall and within-band
  $\rho$ show the proxy holds where the claim lives, not only globally (H5).}
  \label{fig:bridge}
\end{figure}

\paragraph{Weighting robustness.} The deployed overall weights five dimensions
(Eq.~\ref{eq:overall}) on substantive grounds, frozen before this study. To show
the headline does not hinge on the exact (proprietary) coefficients, we recompute
the overall under equal weights and under each leave-one-dimension-out weighting
and report discrimination plus rank agreement with the deployed score
(Table~\ref{tab:weightrobust}, Fig.~\ref{fig:weightrobust}). The equal-weight
score reaches AUROC~\aurocEqualWeightFull and ranks papers nearly identically to
the deployed score ($\rho=\rhoEqualWeight$); every leave-one-out variant stays in
the narrow band AUROC $\in[\looAurocMin,\looAurocMax]$. These are robustness
checks, not a search for a better weighting.

\begin{table}[h]
  \centering
  \caption{Headline discrimination under the deployed, equal, and
  leave-one-dimension-out weightings (full-mini cohort), with rank agreement against
  the deployed score.}
  \label{tab:weightrobust}
  \input{tables/tab_weight_robustness}
\end{table}

\begin{figure}[h]
  \centering
  \includegraphics[width=0.6\columnwidth]{figS_weight_robustness}
  \caption{Reject-vs-accept AUROC (full-mini cohort) under the deployed weighting,
  equal weighting, and each leave-one-dimension-out weighting, with 95\% CIs. The
  result is stable across weightings.}
  \label{fig:weightrobust}
\end{figure}

\paragraph{Temporal leakage controls.} The primary cohort's decisions postdate
the grading model's knowledge cutoff, so the outcome cannot have been memorized;
we bound the residual risks with two comparisons of the full-mini discrimination
(Table~\ref{tab:contam}, Fig.~\ref{fig:contamination}). First, the contaminated
contrast: the fully pre-cutoff ICLR~2025 cohort (AUROC~\aurocRepFull) is no
stronger than the clean primary cohort, so memorization of the outcome is not
inflating the relationship. Second, excluding the \nArxivPrior{} submissions
whose arXiv preprint predates the cutoff leaves the headline essentially
unchanged (AUROC~\aurocNoArxivPriorFull), bounding residual paper-text leakage
(the model reading the manuscript, as a reviewer would, is not itself a confound).
The prospective check (grading a future venue before decisions are released)
remains the only complete control and is reported as follow-up. Whether the
model's priors on authors and institutions move the score is a separate question,
addressed by the prestige-perturbation experiment rather than by a grading
configuration.

\begin{table}[h]
  \centering
  \caption{Temporal leakage controls (full-mini cohort): the clean primary cohort, the
  fully pre-cutoff replication venue (contaminated contrast), and the primary
  cohort with pre-cutoff arXiv papers excluded.}
  \label{tab:contam}
  \input{tables/tab_contamination}
\end{table}

\begin{figure}[h]
  \centering
  \includegraphics[width=0.6\columnwidth]{figS_contamination}
  \caption{Reject-vs-accept AUROC for the clean primary cohort, the pre-cutoff
  replication venue (contaminated contrast), and the primary cohort excluding
  pre-cutoff arXiv papers, with 95\% CIs. The clean cohort is no weaker than the
  contaminated one and is unchanged by the arXiv exclusion: memorization and
  paper-text leakage are not driving the result.}
  \label{fig:contamination}
\end{figure}

\paragraph{Manuscript-length confounding.} A score that simply rewarded longer or
more polished manuscripts would be confounded with length rather than quality. We
correlate the AIPR overall with manuscript page count, word count, reference
count, and figure count on the full-mini cohort (Table~\ref{tab:lengthconfound}). The
rank correlations are weak ($\rho=\lenRhoPage$ with page count,
$\rho=\lenRhoWord$ with word count, $\rho=\lenRhoRefs$ with reference count), so
the discrimination is not explained by manuscript length.

\begin{table}[h]
  \centering
  \caption{Rank correlation of the AIPR overall with manuscript-length metrics
  (full-mini cohort). Weak throughout: the score is not a length proxy.}
  \label{tab:lengthconfound}
  \input{tables/tab_length_confound}
\end{table}

\paragraph{Grading cost.} The cost design (\S\ref{sec:methods}) rests on the
cheap-model full-mini configuration being far cheaper than the frontier full
configuration. Mean tokens used per grading (Table~\ref{tab:cost},
Fig.~\ref{fig:cost}) are \costFullMiniTok{} for full-mini and \costFullTok{} for
full --- the frontier configuration uses \costFullVsMini$\times$ the tokens of
full-mini (the same two-pass pipeline, but the costlier model tier). This is why
the relationship is established on the large full-mini cohort and only confirmed on
the budgeted frontier cohort.

\begin{table}[h]
  \centering
  \caption{Mean tokens used per grading configuration (primary venue). Input =
  prompt/manuscript tokens read; output = generated tokens.}
  \label{tab:cost}
  \input{tables/tab_cost}
\end{table}

\begin{figure}[h]
  \centering
  \includegraphics[width=0.6\columnwidth]{figS_cost}
  \caption{Mean tokens used per grading configuration (primary venue), input vs.\
  output. Both pipeline configurations run the two-pass pipeline (reviewer +
  audit), differing only in the model tier; the naive judge is a single prompt on
  the frontier model.}
  \label{fig:cost}
\end{figure}

\section{Naive-judge baseline (V1): full detail}
\label{app:naive}
This appendix is the detailed companion to the value comparison of
\S\ref{sec:res-value} (the pre-registered experiment V1): the verbatim prompt, the
matched-accept-rate operating points, and the full table. The experiment
separates the contribution of the AIPR pipeline from the contribution of the
underlying model: we grade every frontier-cohort submission a second way---the
same model reads the same submitted PDF, but in place of the two-pass
rubric-and-audit pipeline it receives a single one-paragraph prompt and returns
one overall score with an accept/reject call. We publish the prompt verbatim, so
the baseline is realistic rather than a strawman or a covertly optimized variant:

\begin{quote}
\emph{I'm reviewing this paper for ICLR, a top machine learning conference. Read
it and tell me how good it is: give a single overall quality score from 0 to 100
(where higher means more likely to deserve acceptance) and a one or two sentence
reason for the score.}
\end{quote}

Because the naive judge produces only an overall score, the comparison is on
score and decision alone (Table~\ref{tab:naive}, Fig.~\ref{fig:naive}):
discrimination (AUROC), accept/reject agreement at the pre-registered AIPR@60
operating point (predict accept iff overall $\ge 60$), and run-to-run reliability
(median within-paper SD over repeated runs). So the win is not an artifact of
forcing the baseline onto AIPR's cutoff, we also score each method at its own
threshold matched to the human accept rate. The full pipeline leads on all three:
AUROC \aurocFullFull{} against \aurocNaiveFull{}, balanced accuracy at AIPR@60
\opSixtyFullAcc\% against \opSixtyNaiveAcc\%, and median within-paper SD
\fullRunSD{} against \naiveRunSD{} points. The difference is the discrimination
and reliability the engineered pipeline adds over the same model used directly.

\begin{table}[h]
  \centering
  \caption{Naive one-paragraph judge vs.\ the full AIPR pipeline on the frontier
  cohort (same model, same PDF): discrimination (AUROC), accept/reject agreement
  at AIPR@60 (balanced accuracy), and run-to-run reliability (median within-paper
  SD).}
  \label{tab:naive}
  \input{tables/tab_naive_baseline}
\end{table}

\section{Secondary analysis: comment-quality comparison (ReviewBench axis)}
\label{app:reviewbench}
A concurrent benchmark, ReviewBench \citep{reviewbench2026}, compares the
\emph{review comments} of human reviewers and AI systems by the presence of
structural elements (specification, justification, remedy, anchor) and an
LLM-judged ``consequential'' flag. By construction it does not assess whether a
comment is \emph{correct}: its mapping prompt instructs the judge to ``assess
what the reviewer is saying, not whether they are correct.'' We reproduce its
axis and add the axis it omits.

On the ReviewBench axis, AIPR is at parity with other AI systems: all AI sources
saturate the structural metrics, and AIPR's consequential rate
(\aiprConsequential\%) is comparable to R3's (\rThreeConsequential\%). The
difference appears on \emph{correctness} (Table~\ref{tab:commentquality},
Fig.~\ref{fig:commentquality}): AIPR's citation recommendations are grounded in a
real bibliographic record \aiprGroundedness\% of the time (vs.\
\rThreeGroundedness\% for R3) because they are produced by a retrieval-tool audit
rather than free generation, its comment anchors are verbatim-verifiable
\aiprAnchorFidelity\% of the time, and the share of its consequential comments
that are actually correct (\aiprCorrectness\%) exceeds R3's
(\rThreeCorrectness\%) even though R3 raises consequential-\emph{sounding}
comments at a higher rate. A critical comment that is not correct is the failure a
triage signal must avoid, and a presence/intent benchmark cannot see it; this is
the comment-level analogue of the score-level claim.

\emph{Status:} the numbers here are computed on synthetic comment data for the
draft; the real comparison runs AIPR as an added source through the (vendored,
self-hosted, judge-swapped) ReviewBench framework. This is a secondary analysis,
not primary evidence; correctness is spot-checked by a judge from a different
model family than any system under test.

\begin{table}[h]
  \centering
  \caption{Comment-quality comparison. AI systems are at parity on the
  ReviewBench (form + intent) axis; AIPR leads on the correctness axis ReviewBench
  omits.}
  \label{tab:commentquality}
  \input{tables/tab_comment_quality}
\end{table}

\begin{figure}[h]
  \centering
  \includegraphics[width=\textwidth]{figS_comment_quality}
  \caption{Left: ReviewBench's form+intent axis (structure, consequential) ---
  near parity across AI systems. Right: the correctness axis ReviewBench omits
  (anchor fidelity, citation groundedness, correctness of consequential comments)
  --- AIPR leads, with 95\% Wilson intervals.}
  \label{fig:commentquality}
\end{figure}

\section{Failure-mode case studies}
\label{app:cases}
% PLACEHOLDER: matched table of AIPR weaknesses / omitted citations vs. verbatim
% human-review excerpts for 6--10 examples spanning the patterns in
% Section~\ref{sec:failure}.
\begin{table}[h]
  \centering
  \caption{Matched case studies: AIPR-flagged weakness vs.\ human-review excerpt.
  (Populated from the real export; synthetic placeholder shown in drafts.)}
  \label{tab:casestudy}
  \begin{tabular}{p{0.18\columnwidth}p{0.36\columnwidth}p{0.36\columnwidth}}
    \toprule
    Submission & AIPR weakness / omitted citation & Human-review excerpt \\
    \midrule
    --- & --- & --- \\
    \bottomrule
  \end{tabular}
\end{table}

\section{Additional threats to validity}
\label{app:limits}
The main-text limitations (\S\ref{sec:limits}) cover the load-bearing threats.
The additional, lower-order caveats below are recorded here for completeness.

\paragraph{Exclusion (arXiv-twin) bias.}
We exclude submissions whose self-identity step flags them as an already-published
manuscript under a different title, removing an arXiv-twin leakage risk for the
citation audit (\S\ref{sec:data}). Excluding papers on a content-derived signal
could in principle bias the retained sample. We therefore report the headline
discrimination and rating correlation both with and without the exclusion, and
report the excluded set's outcome distribution; the signal of interest is that the
two agree---i.e.\ the exclusion does not manufacture the effect---and the excluded
count is small relative to the cohort.

\paragraph{Submitted vs.\ post-rebuttal.}
We grade the submitted PDF, but reviewer ratings can shift after rebuttal and
revision. The score is computed on the artifact reviewers first saw; any residual
mismatch between that artifact and the final rating adds noise against us rather
than in our favor.

\paragraph{Grounded is not relevant.}
The citation-grounding advantage (Appendix~\ref{app:reviewbench}) establishes that
a recommended reference is a real record returned by the retrieval tool, not that
it is the \emph{most relevant} omission. Groundedness rules out hallucinated
citations; it does not certify editorial judgment about which prior work matters
most, which remains a human call.

\paragraph{Run-to-run variance.}
The score is stochastic across runs (median within-paper SD
$\approx\runSDmedian$ points on the frontier configuration). All point estimates
use per-submission means and the reported intervals incorporate sampling variance;
the variance is small relative to the between-tier separation but is not zero.
